\chapter{Direct Preference Optimization and Its Family}\label{ch:dpo}

\Cref{ch:rlhf} built the classical RLHF pipeline: collect human comparisons, fit a reward
model, then run PPO with four networks in memory (policy, reference, reward, value) and a
generation engine in the loop. It works, but it is the most complex training system in this book
outside large-scale pretraining. In 2023, Rafailov and colleagues showed that for the
KL-regularised objective RLHF optimises, the reward-modelling stage and the RL stage can be
\emph{collapsed into one supervised loss on the preference pairs themselves}. The result,
\emph{Direct Preference Optimization} (DPO), needs no reward model, no value function, no sampling
during training and no PPO machinery. A DPO step is two forward passes over (prompt, response)
pairs and one backward pass. That is the cost profile of supervised fine-tuning, which makes DPO
easy to run, to scale and to debug.

For a systems person this chapter has two halves. The first derives DPO carefully, because
every design choice in an implementation follows from where the terms in the loss come from:
summing over response tokens only, caching the reference, interpreting the logged metrics. The
second half is about running it: data formats, the pair-batch layout, the reference-model cache,
LoRA-DPO, launch recipes, what healthy logs look like, the failure modes, the large family of
variants (cDPO, IPO, KTO, SimPO, ORPO, CPO, online DPO), and what a DPO step costs in memory
and FLOPs.

\begin{figure}[t]
\centering
\resizebox{\linewidth}{!}{%
\begin{tikzpicture}[
  st/.style={draw, rounded corners=3pt, minimum height=1.0cm, text width=2.35cm, align=center,
    font=\small, line width=0.6pt},
  th/.style={st, fill=fcompute!15, draw=sysblue},
  im/.style={st, fill=fdata!15, draw=keygreen},
  sy/.style={st, fill=fmemory!18, draw=napkin},
  num/.style={circle, fill=labpurple, text=white, font=\scriptsize\bfseries, inner sep=1.2pt},
  arr/.style={-{Stealth[length=2mm]}, line width=0.7pt, draw=black!65}]
\node[th] (a) at (0,0) {KL-regularised RL objective};
\node[th] (b) at (2.9,0) {closed-form optimal policy};
\node[th] (c) at (5.8,0) {implicit reward $\beta\log\frac{\pi}{\pi_{\mathrm{ref}}}$};
\node[th] (d) at (8.7,0) {Bradley--Terry $\Rightarrow$ DPO loss};
\node[im] (e) at (8.7,-1.7) {pair batches, masks, log-probs};
\node[im] (f) at (5.8,-1.7) {reference cache, LoRA-DPO};
\node[im] (g) at (2.9,-1.7) {training recipe \& metrics};
\node[im] (h) at (0,-1.7) {failure modes};
\node[sy] (i) at (0,-3.4) {the family: IPO, KTO, SimPO, ORPO, \ldots};
\node[sy] (j) at (2.9,-3.4) {DPO vs PPO vs GRPO};
\node[sy] (k) at (5.8,-3.4) {memory \& FLOPs per step};
\node[sy] (l) at (8.7,-3.4) {Lab 26: toy task on CPU};
\foreach \x/\n in {a/1,b/2,c/3,d/4,e/5,f/6,g/7,h/8,i/9,j/10,k/11,l/12}
  \node[num] at (\x.north west) {\n};
\draw[arr] (a)--(b); \draw[arr] (b)--(c); \draw[arr] (c)--(d); \draw[arr] (d)--(e);
\draw[arr] (e)--(f); \draw[arr] (f)--(g); \draw[arr] (g)--(h); \draw[arr] (h)--(i);
\draw[arr] (i)--(j); \draw[arr] (j)--(k); \draw[arr] (k)--(l);
\node[font=\scriptsize\itshape, text=sysblue, anchor=east] at (-1.4,0) {theory};
\node[font=\scriptsize\itshape, text=keygreen, anchor=east] at (-1.4,-1.7) {implementation};
\node[font=\scriptsize\itshape, text=napkin, anchor=east] at (-1.4,-3.4) {landscape \& systems};
\end{tikzpicture}}
\caption{Roadmap of this chapter. Steps 1--4 derive DPO; steps 5--8 implement and run it;
steps 9--12 place it among its relatives and cost it out.}
\label{fig:dpo:roadmap}
\end{figure}

% ======================================================================
\section{Why skip the reward model?}\label{sec:dpo:why}

Recall the three-stage pipeline of \cref{ch:rlhf} (\cref{fig:dpo:pipelines}, top). After
supervised fine-tuning (SFT) produces $\pi_{\text{SFT}}$, annotators compare pairs of responses
$(y_w, y_l)$ to a prompt $x$, with $y_w$ the \emph{winner} (``chosen'') and $y_l$ the \emph{loser}
(``rejected''). A reward model $r_\phi$ is fitted with the Bradley--Terry (BT) likelihood
\begin{equation}
\mathcal{L}_{\text{RM}}(\phi) = -\E_{(x,y_w,y_l)}\big[\log\sigma\big(r_\phi(x,y_w)-r_\phi(x,y_l)\big)\big],
\label{eq:dpo:bt}
\end{equation}
and the policy is then trained online against it, with a penalty for drifting from the reference
$\piref=\pi_{\text{SFT}}$:
\begin{equation}
\max_{\pi_\theta}\ \E_{x\sim\mathcal{D},\,y\sim\pi_\theta(\cdot\mid x)}\big[r_\phi(x,y)\big]
 \;-\;\beta\,\KL\big(\pi_\theta(\cdot\mid x)\,\|\,\piref(\cdot\mid x)\big).
\label{eq:dpo:rlobj}
\end{equation}
The RL stage is where the systems pain lives. Every step needs fresh samples from the current policy,
so a generation engine has to sit in the loop and receive weight updates (\cref{ch:rlsys}). A
critic as large as the policy has to be trained, and the reward model must stay in memory.
PPO's many hyperparameters (clip range, GAE $\lambda$, value-loss coefficient, KL
controller) all interact.

DPO keeps the \emph{objective}~\eqref{eq:dpo:rlobj} and the \emph{data} (the same comparison
pairs) but removes both the explicit reward model and the RL loop
(\cref{fig:dpo:pipelines}, bottom). The trick is that \eqref{eq:dpo:rlobj} has a closed-form
solution, and that solution lets us express the reward in terms of the policy.

\begin{figure}[t]
\centering
\resizebox{\linewidth}{!}{%
\begin{tikzpicture}[
  bx/.style={draw, rounded corners=3pt, minimum height=0.95cm, text width=2.1cm, align=center,
    font=\small, line width=0.6pt},
  data/.style={bx, fill=fdata!15, draw=keygreen},
  model/.style={bx, fill=fcompute!15, draw=sysblue},
  loop/.style={bx, fill=fcomm!14, draw=labpurple},
  arr/.style={-{Stealth[length=2mm]}, line width=0.7pt, draw=black!65}]
\node[font=\bfseries\small, anchor=west] at (-1.2,1.15) {RLHF (Chapter 25)};
\node[data] (p1) at (0,0) {preference pairs $(x,y_w,y_l)$};
\node[model] (rm) at (3.0,0) {fit reward model $r_\phi$ (BT loss)};
\node[loop] (ppo) at (6.3,0) {PPO: sample $y\sim\pi_\theta$, score, GAE, update};
\node[model] (out1) at (9.6,0) {aligned $\pi_\theta$};
\draw[arr] (p1)--(rm); \draw[arr] (rm)--(ppo); \draw[arr] (ppo)--(out1);
\node[font=\scriptsize, align=center, text=black!60] at (6.3,-1.2) {online loop: policy + reference + reward + value\\ + generation engine in memory};
\node[font=\bfseries\small, anchor=west] at (-1.2,-2.0) {DPO (this chapter)};
\node[data] (p2) at (0,-3.1) {preference pairs $(x,y_w,y_l)$};
\node[model] (ref) at (3.0,-3.1) {reference log-probs (computed once)};
\node[loop, fill=fdata!10] (dpo) at (6.3,-3.1) {supervised loss on pairs: 2 forwards + 1 backward};
\node[model] (out2) at (9.6,-3.1) {aligned $\pi_\theta$};
\draw[arr] (p2)--(ref); \draw[arr] (ref)--(dpo); \draw[arr] (dpo)--(out2);
\node[font=\scriptsize, align=center, text=black!60] at (6.3,-4.25) {policy only (reference cached or\\ = LoRA base); no sampling};
\end{tikzpicture}}
\caption{RLHF and DPO optimise the same KL-regularised objective from the same data. DPO
replaces the reward model and the online PPO loop with a classification-style loss whose
cost profile is that of supervised fine-tuning.}
\label{fig:dpo:pipelines}
\end{figure}

% ======================================================================
\section{Deriving DPO}\label{sec:dpo:derive}

\subsection{The optimal policy of the KL-regularised objective}

Fix a prompt $x$ and drop it from the notation. For a reward $r(y)$, we want the distribution
$\pi$ over responses that maximises $\E_{\pi}[r(y)] - \beta\,\KL(\pi\,\|\,\piref)$.

\begin{proposition}[Gibbs policy]\label{prop:dpo:gibbs}
The unique maximiser is
\begin{equation}
\pi^*(y) = \frac{1}{Z}\,\piref(y)\exp\!\Big(\frac{r(y)}{\beta}\Big),\qquad
Z = \sum_{y}\piref(y)\exp\!\Big(\frac{r(y)}{\beta}\Big).
\label{eq:dpo:gibbs}
\end{equation}
\end{proposition}
\begin{proof}
Write the objective as
$J(\pi)=\sum_y \pi(y)\big[r(y)-\beta\log\frac{\pi(y)}{\piref(y)}\big]$. Substituting
$r(y)=\beta\log\frac{\pi^*(y)}{\piref(y)}+\beta\log Z$ from \eqref{eq:dpo:gibbs} gives
\[
J(\pi)=\sum_y\pi(y)\Big[\beta\log\frac{\pi^*(y)}{\pi(y)}\Big]+\beta\log Z
      =-\beta\,\KL(\pi\,\|\,\pi^*)+\beta\log Z .
\]
$\beta\log Z$ does not depend on $\pi$, and the KL divergence is non-negative and zero only at
$\pi=\pi^*$.
\end{proof}

The proof shows something useful in practice too: maximising the RLHF objective is the same as
minimising the reverse KL to an exponentially tilted reference. Small $\beta$ means a sharp tilt
(the policy chases reward); large $\beta$ keeps $\pi^*$ close to $\piref$.

\subsection{Inverting for the reward}

\Cref{eq:dpo:gibbs} expresses the policy in terms of the reward. Take logs and solve for the
reward instead (restoring $x$):
\begin{equation}
r(x,y) = \beta\log\frac{\pi^*(y\mid x)}{\piref(y\mid x)} + \beta\log Z(x).
\label{eq:dpo:reparam}
\end{equation}
Any reward function corresponds to some optimal policy, and \eqref{eq:dpo:reparam} turns that
around: \emph{the log-ratio of a policy to the reference, scaled by $\beta$, is a reward for which
that policy is optimal}, up to a prompt-dependent constant $\beta\log Z(x)$. Computing $Z(x)$
exactly is hopeless, since it sums over all possible responses. But BT only ever looks at reward
\emph{differences} between two responses to the same prompt, and in that difference $Z(x)$
cancels.

\subsection{Plugging into Bradley--Terry}

Parameterise the policy as $\pi_\theta$ and substitute \eqref{eq:dpo:reparam} into the BT
likelihood \eqref{eq:dpo:bt}:
\begin{equation}
\boxed{\;
\mathcal{L}_{\text{DPO}}(\theta) = -\E_{(x,y_w,y_l)\sim\mathcal{D}}\Big[\log\sigma\Big(
\beta\log\frac{\pi_\theta(y_w\mid x)}{\piref(y_w\mid x)}
-\beta\log\frac{\pi_\theta(y_l\mid x)}{\piref(y_l\mid x)}\Big)\Big].\;}
\label{eq:dpo:loss}
\end{equation}
That is the whole method. Define the \emph{implicit reward}
$\hat r_\theta(x,y)=\beta\log\frac{\pi_\theta(y\mid x)}{\piref(y\mid x)}$ and the margin
$h=\hat r_\theta(x,y_w)-\hat r_\theta(x,y_l)$. Then $\mathcal{L}_{\text{DPO}}=-\E[\log\sigma(h)]$:
DPO is the BT reward-model loss, with the reward model \emph{being} the policy's log-ratio. Fitting the
implicit reward to the preferences moves the policy directly to the optimum of
\eqref{eq:dpo:rlobj} for that reward. This is the sense of the DPO paper's subtitle, ``your
language model is secretly a reward model''.

\begin{keyidea}[What DPO assumes]
DPO inherits two modelling assumptions. (1) Preferences follow Bradley--Terry: the probability that
$y_w\succ y_l$ is a logistic function of a reward difference. (2) We want the optimum of the
\emph{KL-regularised} objective with reference $\piref$. The optimisation is offline: the expectation in
\eqref{eq:dpo:loss} is over a fixed dataset, not over samples from $\pi_\theta$. That last point
is the source of both DPO's simplicity and most of its failure modes (\cref{sec:dpo:failures}).
\end{keyidea}

\subsection{How DPO uses Bradley--Terry implicitly}

Compare DPO with explicit reward modelling term by term. A reward model is an LM trunk with a
scalar head (\cref{ch:rlhf}), trained by \eqref{eq:dpo:bt}, after which the head is thrown away
and only its scores are used in PPO. DPO uses the same likelihood but with a different hypothesis
class for $r$: functions of the form $\beta\log\pi_\theta/\piref$. Two consequences follow.

\begin{itemize}
\item \textbf{The reward is only defined up to $\beta\log Z(x)$.} Implicit rewards of responses to
  the \emph{same} prompt are comparable; rewards across \emph{different} prompts are not. Do not
  rank prompts by implicit reward, and do not use the DPO model as a general-purpose reward model
  for best-of-$N$ across prompts without care.
\item \textbf{The implicit reward is a sum over tokens.} Because
  $\log\pi(y\mid x)=\sum_t\log\pi(y_t\mid x,y_{<t})$, the implicit reward decomposes into per-token
  log-ratios. DPO can therefore be read as learning token-level credit assignment, and later work
  interprets the per-token log-ratios as an optimal advantage-like quantity (Rafailov et al., 2024).
  In practice, plotting per-token log-ratios over a response is an excellent debugging tool:
  you can see \emph{which} tokens the update promoted or suppressed.
\end{itemize}

\subsection{The gradient: a weighted likelihood update}\label{sec:dpo:grad}

Differentiate \eqref{eq:dpo:loss}. With $\frac{d}{dh}[-\log\sigma(h)]=-\sigma(-h)$,
\begin{equation}
\nabla_\theta\mathcal{L}_{\text{DPO}} = -\beta\,\E\Big[\underbrace{\sigma\big(\hat r_\theta(x,y_l)-\hat r_\theta(x,y_w)\big)}_{\text{weight }w\in(0,1)}
\Big(\underbrace{\nabla_\theta\log\pi_\theta(y_w\mid x)}_{\text{raise winner}}
-\underbrace{\nabla_\theta\log\pi_\theta(y_l\mid x)}_{\text{lower loser}}\Big)\Big].
\label{eq:dpo:grad}
\end{equation}
Each pair contributes ``maximum likelihood on the winner, unlikelihood on the loser'', weighted
by how \emph{wrong} the current implicit reward is about that pair. When the implicit reward
already ranks $y_w$ far above $y_l$ the weight goes to zero, and the pair stops contributing
(\cref{fig:dpo:grad}). The weight is what distinguishes DPO from naive ``unlikelihood'' training,
which keeps pushing the loser's probability down forever and degenerates; the original DPO paper
reports exactly this degeneration for the unweighted version.

\begin{figure}[t]
\centering
\begin{tikzpicture}
\begin{axis}[width=0.62\linewidth, height=5.0cm, xmin=-6, xmax=6, ymin=0, ymax=1.05,
  xlabel={margin $h=\hat r(x,y_w)-\hat r(x,y_l)$}, ylabel={per-pair weight},
  grid=major, grid style={black!10}, legend style={font=\scriptsize, at={(0.98,0.97)}, anchor=north east},
  tick label style={font=\scriptsize}, label style={font=\small}]
\addplot[sysblue, very thick, domain=-6:6, samples=100] {1/(1+exp(x))};
\addlegendentry{DPO: $\sigma(-h)$}
\addplot[warnred, very thick, dashed, domain=-6:6, samples=2] {1};
\addlegendentry{unweighted (unlikelihood)}
\draw[black!50, dashed] (axis cs:0,0) -- (axis cs:0,1.05);
\node[font=\scriptsize, align=center] at (axis cs:-3.6,0.55) {pair mis-ranked:\\ full update};
\node[font=\scriptsize, align=center] at (axis cs:3.6,0.32) {pair already\\ ranked: tiny update};
\end{axis}
\end{tikzpicture}
\caption{The DPO gradient weight $\sigma(-h)$ as a function of the implicit-reward margin
(analytic). Mis-ranked pairs get a full likelihood/unlikelihood update; well-ranked pairs fade
out. At initialisation ($\pi_\theta=\piref$) every margin is $0$ and every weight is $1/2$.}
\label{fig:dpo:grad}
\end{figure}

\begin{pitfall}[The prompt cancels, so truncate it freely but consistently]
$\log\pi(x,y_w)-\log\pi(x,y_l)=\log\pi(y_w\mid x)-\log\pi(y_l\mid x)$: the prompt's own
log-probability cancels inside each ratio-difference, both for $\pi_\theta$ and for $\piref$. So you
may sum log-probs over the whole concatenated sequence and get the same loss. But you \emph{must}
then use exactly the same prompt tokens for winner, loser, policy and reference. If the reference
cache was built with one chat template or truncation policy and training uses another, the
prompt terms no longer cancel and the loss is silently wrong. Summing over response tokens only
(as we do) makes the code robust to this.
\end{pitfall}

% ======================================================================
\section{Preference data in practice}\label{sec:dpo:data}

A DPO dataset is a list of triples. On disk, the most portable format is JSON Lines with one
record per comparison:
\begin{lstlisting}[style=plain]
{"prompt": "How do I boil an egg?", "chosen": "Place the egg in ...", "rejected": "Eggs are ...",
 "source": "helpful", "chosen_score": 8, "rejected_score": 3}
\end{lstlisting}
Real datasets arrive in other shapes, and the first job is a converter into this one.
\begin{itemize}
\item \textbf{Paired conversations.} Some corpora store two complete transcripts that share a
  prefix and differ in the final assistant turn, for example Anthropic's HH-RLHF, whose records hold a
  \texttt{chosen} and a \texttt{rejected} dialogue written as alternating
  ``\texttt{Human:}''/``\texttt{Assistant:}'' turns. The prompt is everything up to the last assistant
  turn; the responses are the two final turns. Decide explicitly what to do with multi-turn
  dialogues whose earlier turns differ between the two sides: dropping them is the simplest correct
  choice.
\item \textbf{Rated completions.} Datasets such as UltraFeedback attach several completions with
  scalar ratings to each prompt. Turn them into pairs by taking the best against a random
  lower-rated one (or the worst), drop ties, and optionally keep the score gap as metadata to filter
  or weight pairs.
\item \textbf{Self-generated (on-policy) pairs.} Sample $k$ responses from $\piref$, score them with
  a reward model, a verifier (\cref{ch:eval}) or an LLM judge, and keep (best, worst). Data drawn
  from the reference policy itself matches DPO's assumptions better than responses written by other
  models (\cref{sec:dpo:failures}).
\end{itemize}

\begin{lstlisting}[caption={Loading JSONL pairs and tokenising them with a chat template. Prompts are truncated from the left (keep the latest turns), responses from the right; EOS is appended to both responses.},label={lst:dpo:tokenize}]
def tokenize_pair(rec, encode, eos_id, template=lambda p: p,
                  max_prompt=512, max_response=512):
    if rec["chosen"] == rec["rejected"]:
        return None                        # no preference signal
    p = encode(template(rec["prompt"]))[-max_prompt:]
    c = encode(rec["chosen"])[: max_response - 1] + [eos_id]
    r = encode(rec["rejected"])[: max_response - 1] + [eos_id]
    return PreferencePair(p, c, r, {k: v for k, v in rec.items()
                                    if k not in ("prompt", "chosen", "rejected")})

# e.g. with Qwen3's ChatML template from Chapter 6:
tok = Qwen3Tokenizer("ckpt/qwen3-0.6b/tokenizer.json")
pairs = [tokenize_pair(r, tok.encode, tok.eos_token_id,
                       template=tok.apply_chat_template)
         for r in load_preference_jsonl("data/prefs.jsonl")]
pairs = [p for p in pairs if p is not None]
\end{lstlisting}

\begin{pitfall}[Append EOS, and use the template the SFT model was trained with]
If responses lack EOS, the policy is never taught where a preferred answer \emph{ends}, and
length drifts. If the template differs from the one used to train $\piref$ (\cref{ch:instruct}),
the reference log-probs are computed on out-of-distribution inputs, the ``KL anchor'' points
somewhere odd, and step-0 metrics look strange (see the checklist in \cref{sec:dpo:recipe}).
\end{pitfall}

\paragraph{Look at the data before training.} For a few dozen random pairs, read the prompt and
both responses and ask: do you agree with the label? Is the chosen response systematically
longer? (Compute the mean length difference; on many public datasets it is clearly positive,
which DPO will happily learn as ``longer is better''.) Are there near-duplicates or identical
pairs? Is the rejected response sometimes truncated mid-sentence (a giveaway artefact the model can
key on)? Ten minutes here saves days of confusing runs.

\subsection{Building pairs on-policy}\label{sec:dpo:onpolicy}

When no suitable dataset exists, or when you want pairs that match the policy's own
distribution, generate them. One round is: sample $k$ responses per prompt from the current model
with temperature, score each with whatever judge you trust, and keep the best and the worst if their
scores differ by enough to be a reliable preference (\cref{lst:dpo:onpolicy}). With a verifier
(\cref{ch:eval}) the scores are exact; with a reward model or LLM judge, the gap threshold filters
out pairs where the judge is guessing.

\begin{lstlisting}[caption={One round of on-policy pair construction (\ident{llmsys.prefopt.build_onpolicy_pairs}, abridged). Generation uses the batched, left-padded, KV-cached path from \cref{ch:serving}.},label={lst:dpo:onpolicy}]
@torch.no_grad()
def build_onpolicy_pairs(model, prompts, score_fn, pad_id, eos_id, k=4,
                         max_new_tokens=64, temperature=1.0, min_gap=0.0):
    flat = [list(p) for p in prompts for _ in range(k)]   # k copies each
    outs = generate_batch(model, flat, max_new_tokens, pad_id=pad_id,
                          eos_id=eos_id, temperature=temperature)
    pairs = []
    for i, p in enumerate(prompts):
        cands = sorted((score_fn(p, outs[i * k + j]), outs[i * k + j])
                       for j in range(k))
        (lo, worst), (hi, best) = cands[0], cands[-1]
        if hi - lo > min_gap and best != worst:
            pairs.append(PreferencePair(list(p), best, worst, {"gap": hi - lo}))
    return pairs
\end{lstlisting}

\begin{systemslens}[Pair generation is an inference job: run it like one]
The $k$ samples of a prompt share their prefix, so a serving engine with prefix caching
(\cref{ch:serving}) prefills the prompt once and decodes $k$ branches. At scale, generation runs
offline on an inference fleet (vLLM, SGLang or the mini engine of \cref{ch:serving}), writes JSONL
shards with the sampling parameters and the generating checkpoint's hash, and a separate scoring job
(a reward model is a prefill-only workload; a judge is another generation job) adds scores. Keep
generation, scoring and training as separate, restartable stages that communicate through files,
just as in the data pipelines of \cref{ch:data}. That makes each stage independently scalable
and debuggable, and lets you re-score without re-generating.
\end{systemslens}

% ======================================================================
\section{Implementing DPO}\label{sec:dpo:impl}

\subsection{The pair batch}

The workhorse is one forward pass over a $(2B,T)$ batch: rows $0..B{-}1$ hold prompt+chosen and
rows $B..2B{-}1$ prompt+rejected, in the same order (\cref{fig:dpo:batch}). A boolean
\texttt{loss\_mask} marks response tokens. Padding goes on the right: under a causal mask real
tokens never attend to later padding, so no attention mask is needed for training (contrast with
left padding for generation, \cref{ch:serving}).

\begin{figure}[t]
\centering
\begin{tikzpicture}[cell/.style={draw=black!40, minimum width=0.52cm, minimum height=0.48cm,
  font=\tiny, inner sep=0pt}, x=0.52cm, y=0.48cm]
\newcommand{\row}[5]{% y, nprompt, nresp, npad, label
  \foreach \i in {1,...,#2} \node[cell, fill=black!10] at (\i,#1) {$x$};
  \pgfmathtruncatemacro{\a}{#2+1}\pgfmathtruncatemacro{\b}{#2+#3}
  \foreach \i in {\a,...,\b} \node[cell, fill=#5] at (\i,#1) {$y$};
  \ifnum#4>0 \pgfmathtruncatemacro{\c}{#2+#3+1}\pgfmathtruncatemacro{\d}{#2+#3+#4}
  \foreach \i in {\c,...,\d} \node[cell, fill=white] at (\i,#1) {\textcolor{black!35}{pad}};\fi}
\row{0}{4}{5}{1}{keygreen!35}
\row{-1}{3}{4}{3}{keygreen!35}
\row{-2}{5}{3}{2}{keygreen!35}
\row{-3.4}{4}{6}{0}{warnred!30}
\row{-4.4}{3}{3}{4}{warnred!30}
\row{-5.4}{5}{4}{1}{warnred!30}
\node[anchor=east, font=\scriptsize] at (0.3,-1) {\begin{tabular}{r}pair 1\\ pair 2\\ pair 3\end{tabular}};
\node[anchor=east, font=\scriptsize] at (0.3,-4.4) {\begin{tabular}{r}pair 1\\ pair 2\\ pair 3\end{tabular}};
\draw[decorate, decoration={brace, amplitude=4pt}, thick] (11,0.4) -- (11,-2.4)
  node[midway, right=5pt, font=\scriptsize, align=left] {rows $0..B{-}1$: prompt + chosen $y_w$};
\draw[decorate, decoration={brace, amplitude=4pt}, thick] (11,-3) -- (11,-5.8)
  node[midway, right=5pt, font=\scriptsize, align=left] {rows $B..2B{-}1$: prompt + rejected $y_l$};
\node[font=\scriptsize, anchor=west] at (0.5,-6.7) {\tikz\node[cell,fill=black!10]{};~prompt (mask 0)\quad
  \tikz\node[cell,fill=keygreen!35]{};/\tikz\node[cell,fill=warnred!30]{};~response incl.\ EOS (mask 1)\quad
  \tikz\node[cell,fill=white]{};~right padding (mask 0)};
\end{tikzpicture}
\caption{The DPO pair batch: one $(2B,T)$ tensor, chosen half first, same pair order in both
halves, right-padded. Log-probabilities are summed only where the mask is 1, giving
$\log\pi(y\mid x)$ for each row; splitting the result at $B$ yields winner and loser terms.}
\label{fig:dpo:batch}
\end{figure}

\begin{lstlisting}[caption={Building the pair batch (from \ident{llmsys/prefopt.py}).},label={lst:dpo:batch}]
def build_pair_batch(pairs, pad_id):
    rows, masks = [], []
    for side in ("chosen", "rejected"):          # chosen half first
        for pr in pairs:
            resp = getattr(pr, side)
            rows.append(list(pr.prompt) + list(resp))
            masks.append([False] * len(pr.prompt) + [True] * len(resp))
    ids, loss_mask = _pad_rows(rows, masks, pad_id)   # right padding
    return {"input_ids": ids, "loss_mask": loss_mask}
\end{lstlisting}

\subsection{Sequence log-probabilities}

\begin{lstlisting}[caption={Response-only sequence log-probabilities. The logit at position $t$ scores token $t{+}1$, so both targets and mask shift by one. \code{gather} minus \code{logsumexp} avoids materialising a $(2B,T,V)$ log-softmax tensor.},label={lst:dpo:logps}]
def sequence_logps(model, input_ids, loss_mask):
    logits = model(input_ids)[:, :-1].float()        # (2B, T-1, V)
    targets = input_ids[:, 1:]
    mask = loss_mask[:, 1:].to(logits.dtype)
    tok = (logits.gather(-1, targets.unsqueeze(-1)).squeeze(-1)
           - torch.logsumexp(logits, dim=-1))         # log p(target)
    return (tok * mask).sum(-1), mask.sum(-1)         # sum, #tokens
\end{lstlisting}

\begin{pitfall}[Sum, don't average, for DPO]
\Cref{eq:dpo:loss} uses $\log\pi(y\mid x)$, the \emph{sum} over response tokens. Averaging per token
silently changes the objective: it removes the dependence on length from both terms and
corresponds to a different (length-normalised) implicit reward. That is a legitimate design,
but it is SimPO's (\cref{sec:dpo:family}), not DPO's, and it changes the scale on which $\beta$
operates. Pick one deliberately and log which one you used.
\end{pitfall}

\subsection{The loss}

\begin{lstlisting}[caption={DPO with optional conservative label smoothing (cDPO, \cref{sec:dpo:family}). The rewards are returned detached for logging.},label={lst:dpo:loss}]
def dpo_loss(pol_c, pol_r, ref_c, ref_r, beta=0.1, label_smoothing=0.0):
    r_c = beta * (pol_c - ref_c)              # implicit reward of y_w
    r_r = beta * (pol_r - ref_r)              # implicit reward of y_l
    h = r_c - r_r
    loss = (-(1 - label_smoothing) * F.logsigmoid(h)
            - label_smoothing * F.logsigmoid(-h))
    return loss, r_c.detach(), r_r.detach()
\end{lstlisting}

Use \code{F.logsigmoid}, never \code{torch.log(torch.sigmoid(h))}: for large negative $h$ the
latter underflows to $\log 0$.

\subsection{The reference model: cache it, or get it for free}\label{sec:dpo:ref}

$\piref$ is frozen, so $\log\piref(y_w\mid x)$ and $\log\piref(y_l\mid x)$ are constants of the
dataset. There are three ways to provide them.
\begin{enumerate}
\item \textbf{A second model on the device.} Simplest to write; costs a full copy of the weights in
  memory and an extra forward per step. A common two-GPU recipe puts the policy on one device and
  the reference on another.
\item \textbf{Precompute and cache} (\cref{lst:dpo:cache}). One forward pass over the dataset before
  training stores two floats per pair. Training then never loads the reference. For multiple
  epochs this also saves compute; for one epoch it costs the same FLOPs, but they are spent
  \emph{before} training, so every GPU byte can go to the policy during training.
\item \textbf{LoRA-DPO.} If the policy is $\piref$ plus LoRA adapters (\cref{sec:instruct:lora}),
  the reference \emph{is} the frozen base: zero the adapters' scaling and run the same model
  (\cref{lst:dpo:lora}). No second copy, no cache file, and the reference stays consistent with the
  policy's exact numerics.
\end{enumerate}

\begin{lstlisting}[caption={Caching reference log-probs (abridged). The optional third column stores KTO's mismatched-pair term (\cref{sec:dpo:family}).},label={lst:dpo:cache}]
@torch.no_grad()
def precompute_reference_logps(ref_model, pairs, pad_id, batch_size=16):
    ref_model.eval()
    out = []
    for s in range(0, len(pairs), batch_size):
        b = build_pair_batch(pairs[s: s + batch_size], pad_id)
        lp, _ = sequence_logps(ref_model, b["input_ids"], b["loss_mask"])
        B = lp.shape[0] // 2
        out.append(torch.stack([lp[:B], lp[B:]], dim=1).float())
    return torch.cat(out)                       # (N, 2): [chosen, rejected]

torch.save({"ref_logps": cache, "ref_ckpt": "ckpt/sft-step-4000",
            "template": "chatml-v1", "n_pairs": len(pairs)}, "cache/ref.pt")
\end{lstlisting}

\begin{lstlisting}[caption={LoRA-DPO: the base weights are the reference model.},label={lst:dpo:lora}]
@contextmanager
def disable_lora(model):
    mods = [m for m in model.modules() if isinstance(m, LoRALinear)]
    saved = [m.scaling for m in mods]
    for m in mods:
        m.scaling = 0.0                   # y = base(x) + 0 * B(A(x))
    try:
        yield model
    finally:
        for m, s in zip(mods, saved):
            m.scaling = s

policy = apply_lora(load_sft_model(), rank=16, alpha=32)
with disable_lora(policy):
    ref_logps = precompute_reference_logps(policy, pairs, pad_id)
\end{lstlisting}

\begin{pitfall}[Cache hygiene]
A reference cache is only valid for the exact (reference checkpoint, tokenizer, template,
truncation lengths, data order) it was built with. Store these alongside the tensor (as in
\cref{lst:dpo:cache}) and assert them at load time. Index the cache by \emph{example id}, not
by batch position, so shuffling cannot misalign it. One cheap end-to-end check: at step 0 the DPO
loss must equal $\ln 2\approx 0.6931$ to three decimals (next section). If it does not, the cache
and the policy disagree.
\end{pitfall}

\subsection{The training step}

\begin{lstlisting}[caption={One optimizer step of DPO with gradient accumulation (the lab's \ident{dpo_train_loop}, abridged).},label={lst:dpo:step}]
for step in range(steps):
    for _ in range(grad_accum):
        idx = next_indices(batch_size)                       # shuffled
        b = build_pair_batch([pairs[i] for i in idx], pad_id)
        lp, _ = sequence_logps(policy, b["input_ids"], b["loss_mask"])
        B = lp.shape[0] // 2
        loss, rc, rr = dpo_loss(lp[:B], lp[B:],
                                ref_logps[idx, 0], ref_logps[idx, 1], beta)
        (loss.mean() / grad_accum).backward()
        log_metrics(loss, rc, rr, lp[:B], lp[B:])
    torch.nn.utils.clip_grad_norm_(policy.parameters(), 1.0)
    opt.step(); opt.zero_grad(set_to_none=True)
\end{lstlisting}

% ======================================================================
\section{Running DPO: a recipe}\label{sec:dpo:recipe}

\subsection{Hyperparameters}

\begin{table}[h]
\centering\small
\begin{tabularx}{\linewidth}{@{}lXl@{}}
\toprule
\textbf{Knob} & \textbf{What it does} & \textbf{Typical starting point}\\
\midrule
$\beta$ & strength of the KL anchor; scale of implicit rewards & $0.1$ (sweep $0.01$--$0.5$)\\
learning rate & full fine-tuning is very sensitive & $\sim 5\times10^{-7}$--$10^{-6}$ (full), $\sim10^{-5}$--$10^{-4}$ (LoRA)\\
effective batch & pairs per optimizer step & 32--128 pairs\\
epochs & passes over fixed data; overfitting comes fast & 1--3\\
warmup & avoids a large early step & a few \% of steps, linear\\
optimizer & the DPO paper used RMSprop; AdamW is common & AdamW, no weight decay\\
max lengths & prompt/response truncation & cover $\geq95\%$ of pairs; log what is cut\\
\bottomrule
\end{tabularx}
\caption{DPO hyperparameters and commonly used starting values. These are conventions, not laws:
the right learning rate in particular depends on the model size and on full vs LoRA training.}
\label{tab:dpo:hparams}
\end{table}

\subsection{Project layout and launch}

\begin{lstlisting}[style=plain]
dpo/
  configs/qwen3-0.6b-dpo.yaml   # model, data, beta, lr, batch, lengths, seeds
  data/prefs.train.jsonl         # {"prompt","chosen","rejected", ...}
  data/prefs.val.jsonl           # ~200-1000 held-out pairs
  cache/ref.train.pt             # reference log-probs + metadata
  runs/<name>/metrics.jsonl      # one line per step
  runs/<name>/ckpt-best/         # best val reward-accuracy checkpoint
\end{lstlisting}

\begin{lstlisting}[style=shell]
# 1) reference cache (once per dataset x reference checkpoint)
python -m dpo.cache_ref --config configs/qwen3-0.6b-dpo.yaml
# 2a) single GPU (LoRA) or one GPU each for policy/reference
python -m dpo.train --config configs/qwen3-0.6b-dpo.yaml
# 2b) full fine-tuning of a larger model: shard the policy (Chapters 14-15)
torchrun --nproc_per_node 8 -m dpo.train --config configs/qwen3-8b-dpo.yaml --fsdp
\end{lstlisting}
The modules named here are the pattern to build, not files in this repository; Lab 26's
\ident{llmsys/prefopt.py} implements their core (\ident{precompute_reference_logps},
\ident{train_preference}). For multi-GPU full fine-tuning, the policy is wrapped with FSDP exactly
as in \cref{ch:sharding}; the cached reference log-probs are just another per-example field
that the data loader returns.

\subsection{What healthy logs look like}

Log, every step: the loss; \textbf{reward accuracy} (fraction of pairs with $\hat r_w>\hat r_l$); the
\textbf{margin} $\E[h]$; the mean implicit rewards of chosen and rejected; the mean
\textbf{log-probs} $\log\pi_\theta(y_w\mid x)$ and $\log\pi_\theta(y_l\mid x)$; and the gradient norm.
\begin{itemize}
\item \textbf{Step 0}: $\pi_\theta=\piref$, so every implicit reward is exactly 0: loss $=\ln 2$,
  margin $=0$, and reward accuracy is $0.5$ if ties count half (0 if you use a strict $>$). Anything
  else is a bug (cache mismatch, template mismatch, dropout left on in one pass, a different dtype for
  reference and policy).
\item \textbf{Early training}: the margin rises, accuracy climbs fast, and the rejected reward goes
  negative much faster than the chosen reward rises.
\item \textbf{Watch the chosen log-prob.} It often \emph{decreases} even though its implicit reward
  increases relative to the loser (\cref{sec:dpo:failures}). A small drop is normal; a collapse
  of both log-probs by tens of nats signals over-optimisation.
\item \textbf{Validation reward accuracy} on held-out pairs is the cheapest model-selection signal.
  Keep the checkpoint with the best value, then confirm with generation-based evaluation.
\end{itemize}

\begin{table}[t]
\centering\small
\begin{tabularx}{\linewidth}{@{}>{\raggedright\arraybackslash}p{3.3cm}>{\raggedright\arraybackslash}X>{\raggedright\arraybackslash}p{4.2cm}@{}}
\toprule
\textbf{Symptom} & \textbf{Likely cause} & \textbf{Check}\\
\midrule
step-0 loss $\neq\ln 2$ & reference cache built with a different template, truncation, dtype or checkpoint; dropout on & recompute 8 pairs on the fly and diff against the cache\\
reward accuracy stuck at 0.5 & loss mask covers prompt only or is shifted wrongly; learning rate far too small & assert mask sums equal response lengths; overfit 16 pairs\\
train accuracy $\to1$ in a few hundred steps, val flat & memorising a small or duplicated set & dedupe, hold out by prompt, fewer epochs, larger $\beta$\\
both log-probs fall by tens of nats & over-optimisation / likelihood displacement & lower lr or step count, raise $\beta$, add NLL on $y_w$\\
generations get longer every checkpoint & chosen responses are longer on average & length stats of the data; LC win rate; SimPO/R-DPO\\
NaN loss & $\log\sigma$ computed as $\log(\sigma)$; fp16 logits & use \code{F.logsigmoid}; fp32 log-softmax\\
OOM at the LM head & $(2B,T,V)$ logits materialised & gather$-$logsumexp, chunking, response slicing\\
\bottomrule
\end{tabularx}
\caption{A DPO debugging checklist.}
\label{tab:dpo:debug}
\end{table}

When something looks wrong, \cref{tab:dpo:debug} lists the usual suspects in the order to check
them.

\subsection{Evaluating a DPO model}

Implicit-reward accuracy measures how well the model \emph{ranks} the dataset's pairs, not how good its
\emph{generations} are. For the latter, generate on held-out prompts and evaluate as in
\cref{ch:eval}: pairwise win rate against the SFT model with an LLM judge
(\cref{sec:eval:judge}), with position swapping to remove order bias, and a length-controlled
variant that regresses out the effect of response length, since judges favour longer answers.
Also run capability benchmarks (multiple-choice knowledge, math with a verifier) and safety probes
(\cref{ch:alignment}) to detect an \emph{alignment tax}, where preference tuning costs capability.

% ======================================================================
\section{DPO on a toy task you can run on a laptop}\label{sec:dpo:toy}

Lab 26's task makes the preferred behaviour \emph{exactly measurable}. A prompt is
\texttt{Q d1 \ldots dk >}; a response is a style token followed by the reversed digits and EOS. The
reference policy is a tiny Qwen3-shaped model (two layers, width 64) trained by SFT on
demonstrations that use the polite style \texttt{sure} only 25\% of the time (otherwise
\texttt{nope}). Preference pairs come in two kinds: \emph{style} pairs (70\%: same correct digits,
\texttt{sure} $\succ$ \texttt{nope}) and \emph{quality} pairs (30\%: both polite, correct reversal
$\succ$ one corrupted digit). Two metrics are computed exactly: $P(\texttt{sure})$, the
probability of the polite first token averaged over prompts, and greedy exact-match accuracy of the
reversal.

\begin{figure}[t]
\centering
\resizebox{\linewidth}{!}{\begin{tikzpicture}
\begin{groupplot}[group style={group size=3 by 1, horizontal sep=1.1cm},
  width=0.385\linewidth, height=5.0cm, xmin=0, xmax=150, grid=major, grid style={black!10},
  tick label style={font=\scriptsize}, label style={font=\scriptsize}, title style={font=\small},
  xlabel={DPO step}]
\nextgroupplot[title={$P(\texttt{sure})$}, ymin=0, ymax=1.05, legend style={font=\tiny, at={(0.97,0.03)}, anchor=south east}]
\addplot[warnred, thick, mark=*, mark size=1.2pt] coordinates {(0,0.239) (25,0.980) (50,0.988) (75,0.994) (100,0.997) (125,0.998) (150,0.999)};
\addlegendentry{$\beta=0.05$}
\addplot[sysblue, thick, mark=*, mark size=1.2pt] coordinates {(0,0.239) (25,0.975) (50,0.968) (75,0.979) (100,0.975) (125,0.984) (150,0.981)};
\addlegendentry{$\beta=0.1$}
\addplot[keygreen, thick, mark=*, mark size=1.2pt] coordinates {(0,0.239) (25,0.962) (50,0.940) (75,0.949) (100,0.954) (125,0.951) (150,0.943)};
\addlegendentry{$\beta=0.5$}
\nextgroupplot[title={mean $\log\pi_\theta(y\mid x)$ (nats)}, ymin=-85, ymax=5, legend columns=2, legend style={font=\tiny, at={(0.5,-0.38)}, anchor=north}]
\addplot[warnred, thick, mark=*, mark size=1.2pt] coordinates {(25,-0.355) (50,-0.697) (75,-2.853) (100,-3.096) (125,-1.975) (150,-5.395)};
\addlegendentry{chosen, $\beta=0.05$}
\addplot[warnred, thick, dashed, mark=o, mark size=1.2pt] coordinates {(25,-19.152) (50,-35.495) (75,-54.346) (100,-61.676) (125,-72.560) (150,-78.341)};
\addlegendentry{rejected, $\beta=0.05$}
\addplot[keygreen, thick, mark=*, mark size=1.2pt] coordinates {(25,-2.606) (50,-2.535) (75,-2.269) (100,-1.578) (125,-1.344) (150,-1.490)};
\addlegendentry{chosen, $\beta=0.5$}
\addplot[keygreen, thick, dashed, mark=o, mark size=1.2pt] coordinates {(25,-11.839) (50,-21.936) (75,-24.506) (100,-25.289) (125,-27.211) (150,-29.750)};
\addlegendentry{rejected, $\beta=0.5$}
\nextgroupplot[title={greedy reversal accuracy}, ymin=0, ymax=1.05, legend style={font=\tiny, at={(0.03,0.03)}, anchor=south west}]
\addplot[warnred, thick, mark=*, mark size=1.2pt] coordinates {(0,1.000) (25,0.940) (50,0.880) (75,0.820) (100,0.980) (125,0.880) (150,0.860)};
\addplot[sysblue, thick, mark=*, mark size=1.2pt] coordinates {(0,1.000) (25,0.960) (50,0.880) (75,0.960) (100,0.660) (125,0.880) (150,0.560)};
\addplot[keygreen, thick, mark=*, mark size=1.2pt] coordinates {(0,1.000) (25,0.940) (50,0.760) (75,0.920) (100,0.960) (125,0.960) (150,0.940)};
\end{groupplot}
\end{tikzpicture}}
\caption{DPO on the toy task, measured on CPU in this environment (one seed; accuracy curves use
50 prompts per point, so they are noisy). Left: the preferred behaviour appears within 25 steps
from the reference's $0.24$. Middle: the margin grows almost entirely by pushing the
\emph{rejected} responses down, by up to 78 nats for $\beta=0.05$; the chosen log-prob also drifts
down. Right: the reversal skill the pairs do not ask about degrades and fluctuates, most for the smaller $\beta$
values: an alignment tax.}
\label{fig:dpo:toycurves}
\end{figure}

\begin{table}[t]
\centering\small
\begin{tabular}{@{}lccl@{}}
\toprule
\textbf{Method (150 steps, lr $10^{-3}$)} & $P(\texttt{sure})$ & \textbf{greedy acc.} & \textbf{note}\\
\midrule
reference (SFT) & 0.239 & 1.00 & starting point\\
DPO, $\beta=0.05$ & 0.998 & 0.85 & \\
DPO, $\beta=0.1$ & 0.981 & 0.52 & \\
DPO, $\beta=0.5$ & 0.947 & 0.95 & stronger anchor\\
cDPO, $\beta=0.1$, $\varepsilon=0.1$ & 0.428 & 0.58 & finite target margin\\
IPO, $\tau=0.1$ & 0.977 & 1.00 & finite target margin\\
KTO, $\beta=0.1$ & 0.999 & 0.98 & pairs split into good/bad\\
SimPO, $\beta=2$, $\gamma=0.5$ & 0.838 & 0.98 & no reference\\
ORPO, $\lambda=0.5$ & 0.999 & 1.00 & no reference, + NLL\\
CPO, $\beta=0.1$ & 0.999 & 0.98 & no reference, + NLL\\
LoRA-DPO, $r=4$, lr $5\times10^{-3}$ & 1.000 & 0.48 & 8516 of 84804 params trained\\
\bottomrule
\end{tabular}
\caption{Final metrics on the toy task (measured on CPU, one seed; accuracy on 100 prompts). Every
method moves the style; methods with an NLL term (ORPO, CPO) or a bounded margin (IPO) keep the
reversal skill best in this run, while plain DPO with small $\beta$ and an aggressive learning rate
pays the largest alignment tax.}
\label{tab:dpo:toyresults}
\end{table}

\begin{figure}[t]
\centering
\begin{tikzpicture}[tok/.style={draw=black!40, minimum width=1.25cm, minimum height=0.62cm, font=\scriptsize\ttfamily, align=center, inner sep=1pt}]
\node[font=\scriptsize, anchor=east] at (-0.9,0.65) {prompt: \texttt{Q 3 8 1 6 >}};
\node[font=\scriptsize, anchor=east] at (-0.9,0.00) {chosen $y_w$};
\node[tok, fill=keygreen!62] at (0.00,0.00) {sure\\[-1pt]{\rmfamily +1.39}};
\node[tok, fill=white] at (1.30,0.00) {6\\[-1pt]{\rmfamily +0.00}};
\node[tok, fill=white] at (2.60,0.00) {1\\[-1pt]{\rmfamily +0.00}};
\node[tok, fill=white] at (3.90,0.00) {8\\[-1pt]{\rmfamily +0.00}};
\node[tok, fill=white] at (5.20,0.00) {3\\[-1pt]{\rmfamily +0.00}};
\node[tok, fill=white] at (6.50,0.00) {\textless eos\textgreater\\[-1pt]{\rmfamily +0.00}};
\node[font=\scriptsize, anchor=east] at (-0.9,-1.00) {style loser};
\node[tok, fill=warnred!53] at (0.00,-1.00) {nope\\[-1pt]{\rmfamily -10.14}};
\node[tok, fill=warnred!12] at (1.30,-1.00) {6\\[-1pt]{\rmfamily -0.97}};
\node[tok, fill=warnred!76] at (2.60,-1.00) {1\\[-1pt]{\rmfamily -15.13}};
\node[tok, fill=warnred!80] at (3.90,-1.00) {8\\[-1pt]{\rmfamily -16.49}};
\node[tok, fill=warnred!80] at (5.20,-1.00) {3\\[-1pt]{\rmfamily -18.31}};
\node[tok, fill=white] at (6.50,-1.00) {\textless eos\textgreater\\[-1pt]{\rmfamily -0.00}};
\node[font=\scriptsize, anchor=east] at (-0.9,-2.00) {quality loser};
\node[tok, fill=keygreen!62] at (0.00,-2.00) {sure\\[-1pt]{\rmfamily +1.39}};
\node[tok, fill=white] at (1.30,-2.00) {6\\[-1pt]{\rmfamily +0.00}};
\node[tok, fill=white] at (2.60,-2.00) {1\\[-1pt]{\rmfamily +0.00}};
\node[tok, fill=warnred!34] at (3.90,-2.00) {9\\[-1pt]{\rmfamily -5.94}};
\node[tok, fill=warnred!66] at (5.20,-2.00) {3\\[-1pt]{\rmfamily -13.07}};
\node[tok, fill=white] at (6.50,-2.00) {\textless eos\textgreater\\[-1pt]{\rmfamily +0.00}};
\end{tikzpicture}
\caption{Per-token log-ratios $\log\pi_\theta(y_t\mid\cdot)-\log\piref(y_t\mid\cdot)$ after 150 DPO
steps ($\beta=0.1$) on the toy task, for one prompt (measured on CPU in this environment; implicit
reward $=\beta\times$ row sum). The winner gains only at the style token ($+1.39\approx\ln(0.99/0.25)$).
In the style loser, the model learned more than ``don't say \texttt{nope}'': every digit
\emph{after} \texttt{nope} became 15--18 nats less likely, mass that went elsewhere. In the quality
loser, the corrupted digit \texttt{9} is penalised, and so is the token after it.}
\label{fig:dpo:tokens}
\end{figure}

\Cref{fig:dpo:toycurves,tab:dpo:toyresults,fig:dpo:tokens} show the three behaviours you will
see at scale. (1) DPO moves exactly the behaviour the pairs encode: $P(\texttt{sure})$ goes from
the reference's $0.24$ to $0.95$--$1.0$ within a few dozen steps. (2) The margin grows mostly by pushing
the \emph{rejected} responses down (tens of nats), while the chosen log-prob barely moves or drifts
down. \Cref{fig:dpo:tokens} shows where: the only token whose probability rises is the style token,
and the digits that follow \texttt{nope} are suppressed far more than the style decision alone
requires. (3) Capability that the pairs do not mention can regress. The reference reverses digits
perfectly, but plain DPO at $\beta=0.1$ and LoRA-DPO with a large learning rate end at about 50\%
greedy accuracy, while the variants with an NLL term (ORPO, CPO) or a bounded margin (IPO) keep it.

cDPO's row is instructive. Its loss is minimised at a finite margin ($\log 9$ in reward units, or 22
nats of log-ratio for $\beta=0.1$), and the cheapest way to reach that margin is to push the losers
down further, \emph{not} to make \texttt{sure} more likely. The objective is satisfied while the
behaviour we care about only moved from $0.24$ to $0.43$. The implicit-reward accuracy and margin
alone would not have shown this, which is why you should always measure the behaviour itself.
With one seed and a toy model, treat the ranking of methods in \cref{tab:dpo:toyresults} as an
illustration of mechanisms, not as a benchmark.

% ======================================================================
\section{Failure modes}\label{sec:dpo:failures}

\paragraph{Likelihood displacement.} \Cref{eq:dpo:grad} only constrains the \emph{difference}
of log-ratios, so both log-probabilities may fall while the margin grows. Probability mass freed
from both responses goes \emph{somewhere else}: sometimes to paraphrases of the winner (harmless),
sometimes to responses unlike either (harmful). Displacement is worse when $y_w$ and $y_l$ are
similar (they share most tokens, so pushing $y_l$ down drags $y_w$ along); Razin et al.\ (2024)
analyse when this happens. Mitigations: an auxiliary SFT/NLL term on the chosen response (as in
iterative RPO, Pang et al., 2024, and in CPO), larger $\beta$, fewer epochs, and on-policy pairs.

\paragraph{Overfitting deterministic preferences.} If a pair is always labelled the same way, the
BT-optimal margin is infinite; DPO keeps pushing $h\to\infty$, so the implicit reward and the
policy can move arbitrarily far from $\piref$ despite the KL term. The KL regularisation in
\eqref{eq:dpo:rlobj} acts through the \emph{optimal policy}, but with finite data the empirical
optimum is not regularised. IPO (\cref{sec:dpo:family}) was designed for exactly this.

\paragraph{Length exploitation.} If chosen responses are longer on average, DPO learns
``longer is better''. A summed log-prob implicit reward also scales with length. Watch the mean
generation length per checkpoint; use length-controlled evaluation; consider length-normalised
variants (SimPO), a length penalty in the margin (R-DPO, Park et al., 2024), or length-balanced data.

\paragraph{Off-policy data.} DPO's gradient only touches the responses in the dataset. If those
were written by a different (stronger) model, the policy is pushed towards sequences it would
rarely sample itself, and what it actually generates is never directly corrected. Generating pairs
from $\piref$ (or iteratively from the current policy; \emph{online DPO}) closes this gap at the
cost of reintroducing a generation engine.

\paragraph{A mis-specified reference.} DPO anchors to whatever $\piref$ is. A reference that is not
the SFT model the data was collected around (wrong checkpoint, wrong template) makes the KL
anchor meaningless. This is why the cache metadata matters.

\begin{systemslens}[Why DPO is easy to scale, and where it isn't]
A DPO step is SFT on $2\times$ the sequences plus (optionally) a no-grad forward of the
reference. Everything from \cref{ch:pretrain,ch:distributed,ch:sharding} applies unchanged:
FSDP, activation checkpointing, mixed precision, sequence packing. There is no generation, no
weight sync to an inference engine, and no long-tailed rollouts (\cref{ch:rlsys}), so throughput is
predictable and GPU utilisation can approach SFT levels. The price is off-policy data. The moment
you go online (sample new pairs every round), you inherit the RL-systems problems of
\cref{ch:rlsys}, but in a milder, batched form: generate a whole dataset, label it, train, repeat.
\end{systemslens}

% ======================================================================
\section{The family: cDPO, IPO, KTO, SimPO, ORPO, CPO and online variants}\label{sec:dpo:family}

DPO kicked off a large literature of ``direct alignment'' losses. Nearly all of them are
functions of the same few per-pair quantities: summed or averaged log-probs of the winner and
loser under the policy, optionally minus the same under a reference. \Cref{fig:dpo:losses} plots
several as functions of the margin, and \cref{tab:dpo:family} summarises them.

\paragraph{Conservative DPO (cDPO) and robust variants.} If a fraction $\varepsilon$ of labels are
flipped, model the label as noisy:
$\mathcal{L}=-(1-\varepsilon)\log\sigma(h)-\varepsilon\log\sigma(-h)$. The optimum is now at a finite
margin, $h^*=\log\frac{1-\varepsilon}{\varepsilon}$, which also tames overfitting.

\paragraph{IPO.} Azar et al.\ (2023) cast RLHF and DPO as instances of a general objective
($\Psi$PO) and showed that DPO's use of the BT logit is what makes deterministic preferences
blow up. \emph{Identity preference optimisation} replaces the logistic loss with a squared
regression of the log-ratio margin onto a finite target:
\begin{equation}
\mathcal{L}_{\text{IPO}}=\E\Big[\Big(\log\frac{\pi_\theta(y_w\mid x)\,\piref(y_l\mid x)}
{\pi_\theta(y_l\mid x)\,\piref(y_w\mid x)}-\frac{1}{2\tau}\Big)^2\Big].
\label{eq:dpo:ipo}
\end{equation}
The margin is pulled \emph{towards} $1/(2\tau)$ and penalised for overshooting it. $\tau$ plays
$\beta$'s role. Many implementations feed length-averaged log-probs into \eqref{eq:dpo:ipo}.

\paragraph{KTO.} Ethayarajh et al.\ (2024) observe that paired comparisons are expensive to
collect, while binary ``this response is good/bad'' signals (thumbs up/down) are cheap. Drawing on
prospect theory, KTO assigns each unpaired example a value relative to a \emph{reference point}
$z_0$:
\begin{align}
r_\theta(x,y)&=\log\frac{\pi_\theta(y\mid x)}{\piref(y\mid x)},\qquad
z_0=\max\!\big(0,\ \widehat{\KL}(\pi_\theta\,\|\,\piref)\big)\ \text{(detached)},\\
v(x,y)&=\begin{cases}\lambda_D\,\sigma\big(\beta(r_\theta(x,y)-z_0)\big) & y\ \text{desirable}\\
\lambda_U\,\sigma\big(\beta(z_0-r_\theta(x,y))\big) & y\ \text{undesirable}\end{cases},\qquad
\mathcal{L}_{\text{KTO}}=\E\big[\lambda_y-v(x,y)\big].
\end{align}
The logistic value function is concave for gains and convex for losses, like Kahneman and
Tversky's (\cref{fig:dpo:kto}). $\lambda_D,\lambda_U$ rebalance classes (setting $\lambda_U>\lambda_D$
encodes loss aversion). $z_0$ is estimated per batch, without sampling, by scoring
\emph{mismatched} pairs (prompt $i$ with the response of example $i{+}1$) under policy and
reference. Ethayarajh et al.\ frame DPO, KTO and others as \emph{human-aware losses} (HALOs).
From a systems view, KTO's per-example (rather than per-pair) batches pack better and allow
unbalanced data.

\paragraph{SimPO.} Meng et al.\ (2024) drop the reference model and use the
\emph{length-normalised} log-probability as the reward, plus a target margin $\gamma$:
$\mathcal{L}=-\log\sigma\big(\frac{\beta}{|y_w|}\log\pi_\theta(y_w\mid x)-\frac{\beta}{|y_l|}\log\pi_\theta(y_l\mid x)-\gamma\big)$.
No reference means no cache and no second model, and the length normalisation counters length
exploitation. The cost: no explicit anchor to $\piref$, so the learning rate and step count
carry the regularisation.

\paragraph{ORPO.} Hong et al.\ (2024) fold SFT and preference learning into one stage with no
reference model: the loss is the SFT negative log-likelihood of $y_w$ plus
$-\lambda\log\sigma\big(\log\frac{\text{odds}_\theta(y_w)}{\text{odds}_\theta(y_l)}\big)$, where
$\text{odds}(y)=p/(1-p)$ with $p$ the length-normalised sequence probability. Starting from a base
model, one ORPO run replaces SFT followed by DPO.

\paragraph{CPO.} Contrastive preference optimisation (Xu et al., 2024), proposed for machine
translation, uses a reference-free DPO-like term $-\log\sigma(\beta(\log\pi_\theta(y_w)-\log\pi_\theta(y_l)))$
plus an NLL term on $y_w$. Like ORPO and SimPO, it trades the KL anchor for simplicity.

\paragraph{Earlier and adjacent methods.} SLiC-HF (Zhao et al., 2023) uses a hinge loss on the
log-probability margin plus a regulariser towards SFT. RSO (Liu et al., 2023) argues for
sampling pairs from (an estimate of) the optimal policy via statistical rejection sampling.
sDPO (Kim et al., 2024) splits the data into stages and uses the previous stage's model as the
next reference. Anchored preference optimisation (APO, D'Oosterlinck et al., 2024) controls whether
the winner's likelihood should go up or down explicitly.

\paragraph{Online and iterative DPO.} Generate $k$ responses per prompt from the \emph{current} policy,
label them (reward model, verifier or judge), form pairs, run a DPO epoch, and repeat. This
recovers much of the benefit of on-policy RL at a fraction of the systems complexity. Self-rewarding
LMs (Yuan et al., 2024) use the model itself as the judge. SPIN (Chen et al., 2024) treats SFT data
as winners against the model's own generations. Nash learning from human feedback (Munos et al.,
2023) treats alignment as a game against a preference model rather than a reward.

\begin{figure}[t]
\centering
\begin{tikzpicture}
\begin{axis}[width=0.7\linewidth, height=5.4cm, xmin=-3, xmax=6, ymin=0, ymax=4.2,
  xlabel={margin $h$ (log-ratio difference, or scaled log-prob difference)}, ylabel={loss},
  grid=major, grid style={black!10}, legend columns=2,
  legend style={font=\scriptsize, at={(0.5,-0.28)}, anchor=north},
  tick label style={font=\scriptsize}, label style={font=\small}]
\addplot[sysblue, very thick, domain=-3:6, samples=120] {ln(1+exp(-x))};
\addlegendentry{DPO $-\log\sigma(h)$}
\addplot[keygreen, very thick, dashed, domain=-3:6, samples=120] {-0.9*ln(1/(1+exp(-x))) - 0.1*ln(1/(1+exp(x)))};
\addlegendentry{cDPO, $\varepsilon=0.1$}
\addplot[warnred, very thick, domain=-3:6, samples=120] {0.25*(x-2.5)^2};
\addlegendentry{IPO $(h-\frac{1}{2\tau})^2$, $\frac{1}{2\tau}=2.5$ (scaled)}
\addplot[napkin, very thick, dotted, domain=-3:6, samples=120] {ln(1+exp(-(x-1)))};
\addlegendentry{SimPO, margin $\gamma=1$}
\addplot[labpurple, thick, domain=-3:6, samples=120] {max(0,1-x)};
\addlegendentry{hinge (SLiC)}
\end{axis}
\end{tikzpicture}
\caption{Direct-alignment losses as functions of the margin (analytic; IPO scaled by 1/4 to
share the axis). DPO never stops rewarding a larger margin; cDPO's minimum is at a finite margin
($\log 9\approx2.2$ for $\varepsilon=0.1$); IPO penalises overshooting its target; SimPO shifts
DPO's curve by the target margin $\gamma$; the hinge stops at margin 1.}
\label{fig:dpo:losses}
\end{figure}

\begin{figure}[t]
\centering
\begin{tikzpicture}
\begin{axis}[width=0.62\linewidth, height=5.0cm, xmin=-8, xmax=8, ymin=-1.6, ymax=1.1,
  xlabel={$r_\theta(x,y)-z_0$}, ylabel={value $v$ (centred)},
  grid=major, grid style={black!10}, legend style={font=\scriptsize, at={(0.02,0.98)}, anchor=north west},
  tick label style={font=\scriptsize}, label style={font=\small}]
\addplot[keygreen, very thick, domain=0:8, samples=60] {1/(1+exp(-0.6*x)) - 0.5};
\addlegendentry{desirable, $\lambda_D=1$}
\addplot[warnred, very thick, domain=-8:0, samples=60] {1.5*(1/(1+exp(0.6*x)) - 0.5)*(-1)};
\addlegendentry{undesirable, $\lambda_U=1.5$}
\draw[black!50] (axis cs:-8,0) -- (axis cs:8,0);
\draw[black!50] (axis cs:0,-1.6) -- (axis cs:0,1.1);
\node[font=\scriptsize, align=center] at (axis cs:4.5,-0.6) {concave for gains};
\node[font=\scriptsize, align=center] at (axis cs:-4.3,0.55) {convex for losses,\\ steeper (loss aversion)};
\end{axis}
\end{tikzpicture}
\caption{The shape of KTO's value function (analytic, $\beta=0.6$): a logistic stand-in for
Kahneman--Tversky's value curve, centred at the reference point $z_0$. Gains (desirable outputs
above the reference point) saturate; losses are weighted by $\lambda_U$.}
\label{fig:dpo:kto}
\end{figure}

\begin{table}[t]
\centering\small
\begin{tabularx}{\linewidth}{@{}lccXl@{}}
\toprule
\textbf{Method} & \textbf{Ref.} & \textbf{Data} & \textbf{Per-example signal} & \textbf{Extra knobs}\\
\midrule
DPO & yes & pairs & $\beta(\Delta\log\pi_\theta-\Delta\log\piref)$, logistic & $\beta$\\
cDPO & yes & pairs & as DPO with label noise $\varepsilon$ & $\beta,\varepsilon$\\
IPO & yes & pairs & squared error to margin $1/(2\tau)$ & $\tau$\\
KTO & yes & unpaired & $\sigma(\pm\beta(r-z_0))$, asymmetric weights & $\beta,\lambda_D,\lambda_U$\\
SimPO & no & pairs & length-normalised log-prob difference, margin $\gamma$ & $\beta,\gamma$\\
ORPO & no & pairs & NLL$(y_w)$ + log odds-ratio term & $\lambda$\\
CPO & no & pairs & ref-free DPO term + NLL$(y_w)$ & $\beta$\\
SLiC-HF & no & pairs & hinge on log-prob margin + SFT regulariser & margin, weight\\
online DPO & yes & pairs (fresh) & DPO on pairs sampled from current policy & sampler, labeller\\
\bottomrule
\end{tabularx}
\caption{The direct-alignment family at a glance. ``Ref.'' is whether a reference model (or cache)
is needed. All except online DPO train on a fixed dataset.}
\label{tab:dpo:family}
\end{table}

\begin{lstlisting}[caption={Four relatives in a few lines each (from \ident{llmsys/prefopt.py}). Sums are over response tokens; \code{len\_*} are token counts.},label={lst:dpo:family}]
def ipo_loss(pc, pr, rc, rr, tau=0.1):
    h = (pc - rc) - (pr - rr)
    return (h - 1.0 / (2.0 * tau)) ** 2

def kto_loss(pol, ref, desirable, z0, beta=0.1, lam_d=1.0, lam_u=1.0):
    r, z0 = pol - ref, z0.detach().clamp(min=0)
    v = torch.where(desirable, lam_d * torch.sigmoid(beta * (r - z0)),
                    lam_u * torch.sigmoid(beta * (z0 - r)))
    lam = torch.where(desirable, torch.full_like(r, lam_d),
                      torch.full_like(r, lam_u))
    return lam - v

def simpo_loss(pc_sum, pr_sum, len_c, len_r, beta=2.0, gamma=0.5):
    h = beta * pc_sum / len_c - beta * pr_sum / len_r
    return -F.logsigmoid(h - gamma)

def orpo_loss(pc_sum, pr_sum, len_c, len_r, lam=0.1):
    a_c, a_r = pc_sum / len_c, pr_sum / len_r          # mean log-probs < 0
    log_odds = lambda a: a - torch.log(-torch.expm1(a))  # log p/(1-p)
    return -a_c - lam * F.logsigmoid(log_odds(a_c) - log_odds(a_r))
\end{lstlisting}

\paragraph{One trainer for every loss.} Because all these losses consume the same per-pair
quantities, one training loop serves them all; only the reduction differs.
\Cref{lst:dpo:dispatch} shows the dispatch in \ident{llmsys.prefopt}. KTO is the only one that
needs an extra forward, over the mismatched pairs used to estimate its reference point. The
reference cache gets a third column for it, so the reference model is still never loaded during
training.

\begin{lstlisting}[caption={Dispatching on the loss type in one training step (\ident{preference_step_loss}, abridged).},label={lst:dpo:dispatch}]
lp, n = sequence_logps(model, b["input_ids"], b["loss_mask"])
pc, pr = split_pair_logps(lp); nc, nr = split_pair_logps(n)
if cfg.needs_reference:                       # dpo, ipo, kto
    rc, rr = ref[:, 0], ref[:, 1]
if cfg.loss == "dpo":
    loss, a, c = dpo_loss(pc, pr, rc, rr, cfg.beta, cfg.label_smoothing)
elif cfg.loss == "ipo":
    loss, a, c = ipo_loss(pc, pr, rc, rr, cfg.ipo_tau)
elif cfg.loss == "simpo":
    loss, a, c = simpo_loss(pc, pr, nc, nr, cfg.beta, cfg.simpo_gamma)
elif cfg.loss == "orpo":
    loss, a, c = orpo_loss(pc, pr, nc, nr, cfg.orpo_lambda)
elif cfg.loss == "kto":                       # winners desirable, losers not
    mism = [mismatched_pair(all_pairs, i) for i in idx]   # prompt i, y_w of i+1
    mb = build_pair_batch(mism, pad_id)
    mlp, _ = sequence_logps(model, mb["input_ids"], mb["loss_mask"])
    z0 = kto_kl_estimate(split_pair_logps(mlp)[0], ref[:, 2])
    ...
metrics = preference_metrics(a, c)
\end{lstlisting}

% ======================================================================
\section{DPO vs PPO vs GRPO}\label{sec:dpo:compare}

\begin{table}[t]
\centering\small
\begin{tabularx}{\linewidth}{@{}lXXX@{}}
\toprule
 & \textbf{PPO-RLHF} (\cref{ch:rlhf}) & \textbf{DPO} (this chapter) & \textbf{GRPO} (\cref{ch:grpo,ch:grpo2})\\
\midrule
data & prompts + reward model & fixed preference pairs & prompts + verifier / reward\\
on/off-policy & on-policy (minibatch reuse) & offline & on-policy (near)\\
reward model & explicit, trained first & implicit $\beta\log\pi_\theta/\piref$ & rule/verifier or RM\\
critic & value network & none & none (group baseline)\\
models in memory & policy, ref, RM, value & policy (+ ref or cache) & policy, ref (optional)\\
generation in loop & yes & no & yes, $G$ samples/prompt\\
step cost & rollouts + 4 forwards + 2 backwards & 2 seq.\ fwd+bwd (+ref fwd) & rollouts + fwd/bwd on $G$ seqs\\
main knobs & clip, GAE $\lambda$, KL coef., value coef. & $\beta$, lr, epochs & $G$, clip, KL coef., lengths\\
strengths & explores, optimises a real reward & simple, stable, cheap & no critic, verifiable rewards\\
weaknesses & complex, expensive, brittle & off-policy, overfits, length bias & needs verifiable reward, rollout cost\\
\bottomrule
\end{tabularx}
\caption{The three post-training workhorses compared. The rows are characteristics of the
methods as usually run; each has online/offline hybrids.}
\label{tab:dpo:compare}
\end{table}

\Cref{tab:dpo:compare} compares the three methods. Which works best is an empirical, setting-dependent question. Some
careful comparisons (e.g.\ Xu et al., 2024, ``Is DPO superior to PPO for LLM alignment?'') report
that a well-tuned PPO can outperform DPO, especially on tasks such as code where on-policy
exploration and an accurate reward matter, and that DPO is sensitive to the distribution shift
between the data and the policy. Others find DPO competitive at much lower cost for chat-style
preference data. A practical reading for a systems person:
\begin{itemize}
\item If you have a \textbf{fixed pairwise dataset} and want a dependable improvement cheaply, start
  with DPO (or SimPO/ORPO if a reference is inconvenient).
\item If you have a \textbf{verifiable reward} (math, code, tests), the on-policy, critic-free
  methods of \cref{ch:grpo} are the natural fit.
\item If you have a \textbf{learned reward model} and the budget for an RL system, PPO or online DPO
  can exploit it, with the reward-hacking caveats of \cref{ch:rlhf}.
\item Many production post-training pipelines use \emph{several}: SFT, then DPO-style
  preference tuning, then RL with verifiable rewards, sometimes iterated.
\end{itemize}

% ======================================================================
\section{Systems: what a DPO step costs}\label{sec:dpo:systems}

\begin{napkinmath}[FLOPs per pair and per epoch]
Take Qwen3-8B ($N\approx8.19\times10^9$ parameters), prompts of 256 tokens and responses of 256
tokens each. A pair is two sequences of 512 tokens: 1{,}024 tokens. The policy costs about $6N$
FLOPs per token (forward + backward), the reference $2N$ (forward only):
\begin{align*}
\text{policy: }& 6\cdot8.19\times10^9\cdot1024\approx5.0\times10^{13}\ \text{FLOPs/pair},\\
\text{reference: }& 2\cdot8.19\times10^9\cdot1024\approx1.7\times10^{13}\ \text{FLOPs/pair}.
\end{align*}
For 160k pairs (about the size of a large public chat preference set) that is $1.07\times10^{19}$
FLOPs per epoch, or about \textbf{7.5 H100-hours} at 40\% MFU
($0.4\times989$\,TFLOP/s). The reference forward is a quarter of the total: caching it moves
$2.7\times10^{18}$ FLOPs into a one-off precomputation that can run on inference-optimised
hardware, and it is saved outright for every additional epoch.
\end{napkinmath}

\begin{napkinmath}[Memory: full DPO vs LoRA-DPO for 8B]
Full fine-tuning with AdamW and fp32 master weights needs $\approx16$ bytes/param
(\cref{ch:pretrain}): $16\times8.19\times10^9\approx131$\,GB of weights, gradients and optimizer
state, plus activations for $2B$ sequences. A bf16 reference adds 16.4\,GB. This does not fit on
one 80\,GB GPU: shard the policy with FSDP/ZeRO-3 (\cref{ch:sharding}) over 8 GPUs ($\approx16$\,GB of
states per GPU) and cache the reference.

LoRA-DPO with rank 16 on every linear layer of Qwen3-8B trains about $44$M parameters: a frozen
bf16 base of 16.4\,GB plus $16\times4.4\times10^7\approx0.7$\,GB of adapter states. The reference
costs nothing extra, because it \emph{is} the base with adapters disabled. Activations then dominate:
use activation checkpointing (\cref{ch:perf}) and keep the logits small (next box).
\end{napkinmath}

\begin{figure}[t]
\centering
\begin{tikzpicture}
\begin{axis}[ybar stacked, width=0.78\linewidth, height=5.6cm, bar width=22pt,
  ylabel={GB (excluding activations)}, ymin=0, ymax=170,
  symbolic x coords={A,B,C,D},
  xtick=data, xticklabels={full + ref model, full + cache, full FSDP/8 (per GPU), LoRA (ref = base)}, x tick label style={font=\scriptsize, text width=2.4cm, align=center},
  legend style={font=\scriptsize, at={(0.5,-0.3)}, anchor=north, legend columns=1},
  tick label style={font=\scriptsize}]
\addplot[fill=fcompute!60, draw=sysblue] coordinates {(A,131) (B,131) (C,16.4) (D,16.4)};
\addplot[fill=fmemory!60, draw=napkin] coordinates {(A,16.4) (B,0) (C,0) (D,0)};
\addplot[fill=fdata!60, draw=keygreen] coordinates {(A,0) (B,0) (C,0) (D,0.7)};
\legend{policy weights+grads+Adam (or frozen base), reference model, LoRA adapter states}
\node[font=\scriptsize, above] at (axis cs:A,147.4) {147.4};
\node[font=\scriptsize, above] at (axis cs:B,131) {131};
\node[font=\scriptsize, above] at (axis cs:C,16.4) {16.4};
\node[font=\scriptsize, above] at (axis cs:D,17.1) {17.1};
\end{axis}
\end{tikzpicture}
\caption{Analytic memory for DPO on an 8B model (16 bytes/param for full training, bf16 frozen
weights). Caching the reference removes 16\,GB; FSDP divides policy state across GPUs; LoRA-DPO
needs only the frozen base plus tiny adapter states, with the reference for free.}
\label{fig:dpo:memory}
\end{figure}

\begin{systemslens}[Logits are the hidden activation hog]
The output layer produces $(2B,T,V)$ logits. For $B=8$ pairs, $T=1024$ and Qwen's $V=151{,}936$
in fp32 that is $2\cdot8\cdot1024\cdot151936\cdot4\approx10$\,GB, before the log-softmax copy.
Three standard fixes: (1) compute log-probs with \code{gather} minus \code{logsumexp}
(\cref{lst:dpo:logps}) and never materialise \code{log\_softmax}; (2) slice to response positions
\emph{before} the LM head when prompts are long; (3) chunk the LM head and log-prob computation
over the sequence (the same trick as chunked cross-entropy in \cref{ch:distill}). Fused
``linear + cross-entropy'' kernels go further by never writing the logits to HBM.
\end{systemslens}

\begin{systemslens}[Share the prompt: one prefix, two continuations]
Winner and loser share the prompt, yet the pair batch processes it twice. With $T_p=256$,
$T_w=T_l=256$, a pair costs $2T_p+T_w+T_l=1024$ tokens; computing the prompt once costs
$T_p+T_w+T_l=768$, saving \textbf{25\%} of FLOPs and activation memory, and much more for long
multi-turn prompts. Two implementations: (a) pack prompt, $y_w$ and $y_l$ into one sequence with a
custom mask in which $y_l$ attends to the prompt but not to $y_w$ (a two-branch tree mask, as in
tree-attention speculative verification, \cref{ch:efficient}), with $y_l$'s positions restarting
after the prompt; or (b) prefill the prompt once and run both continuations against its KV cache
(\cref{ch:serving}), which is natural at inference time and possible at training time with
attention kernels that accept a shared prefix. Either way, verify against the naive pair batch on
a few examples: the log-probs must agree to numerical precision.
\end{systemslens}

\begin{deepdive}[Padding, bucketing and packing pairs]
Pairs have wildly different lengths, so right-padding a random batch to its longest row wastes
FLOPs (\cref{ch:classify}). Bucket pairs by $\max(|x|+|y_w|,|x|+|y_l|)$, or pack several pairs into
fixed-length rows with block-diagonal attention masks and per-segment position ids
(\cref{sec:instruct:packing}). Packing requires keeping track of which segment is the winner and
which the loser of which pair, which is just bookkeeping on the \code{loss\_mask}.
\end{deepdive}

\paragraph{Placement recipes.} (1) \emph{Single GPU, small model}: policy and reference both on the
GPU, or LoRA-DPO. (2) \emph{Two GPUs}: policy on GPU~0, reference on GPU~1; compute the reference
log-probs for the next micro-batch while the policy runs backward on the current one. (3)
\emph{Many GPUs}: FSDP-shard the policy and use a precomputed cache built by an inference job
beforehand. (4) \emph{Online DPO}: add the generation engine of \cref{ch:serving} or \cref{ch:rlsys}
on separate GPUs; generation is batched per round, so the simpler synchronous design of
\cref{ch:rlsys} is usually adequate.

% ======================================================================
\section{Summary}\label{sec:dpo:summary}
\begin{itemize}
\item The KL-regularised RLHF objective has a closed-form optimum $\pi^*\propto\piref e^{r/\beta}$.
  Inverting it writes the reward as $\beta\log\pi/\piref$ up to a prompt constant that cancels in
  Bradley--Terry comparisons. Substituting gives the DPO loss, a logistic loss on the difference
  of log-ratios.
\item The DPO gradient is likelihood on winners and unlikelihood on losers, weighted by how
  mis-ranked each pair is under the implicit reward.
\item Implementation is SFT-like: a $(2B,T)$ pair batch with a response mask, summed response
  log-probs, a reference that is cached or (with LoRA) the frozen base, and the sanity check that the
  step-0 loss equals $\ln 2$.
\item Watch reward accuracy, margins, \emph{and} the raw log-probs of both responses: the chosen
  log-prob often falls (likelihood displacement). Over-optimising fixed data costs capability.
\item Variants trade the reference model, pairing and length handling: cDPO and IPO bound the margin,
  KTO uses unpaired good/bad labels, SimPO, ORPO and CPO drop the reference, and online DPO restores
  on-policy data.
\item DPO costs $\approx(6N+2N)$ FLOPs per token of the pair, has no generation in the loop, and
  scales with the same tools as pretraining. Logit memory and duplicated prompts are the main
  inefficiencies to engineer away.
\end{itemize}

% ======================================================================
\section{Exercises}\label{sec:dpo:exercises}

\begin{exercise}[Gibbs policy, continuous check]
For a toy ``response space'' of three outcomes with $\piref=(0.5,0.3,0.2)$ and $r=(0,1,2)$,
compute $\pi^*$ for $\beta\in\{0.1,1,10\}$ and verify numerically that $\pi^*$ maximises
$\E_\pi[r]-\beta\KL(\pi\|\piref)$ by evaluating the objective at random distributions.
\end{exercise}

\begin{exercise}[Z cancels]
Show that adding any function $f(x)$ to the reward leaves the BT likelihood unchanged and does not
change $\pi^*$ in \eqref{eq:dpo:gibbs}. What does this imply about comparing implicit rewards across
prompts?
\end{exercise}

\begin{exercise}[The gradient]
Derive \eqref{eq:dpo:grad} from \eqref{eq:dpo:loss}. Then show that at initialisation
($\pi_\theta=\piref$) the DPO gradient equals $\beta/2$ times the gradient of
$\log\pi_\theta(y_w\mid x)-\log\pi_\theta(y_l\mid x)$.
\end{exercise}

\begin{exercise}[cDPO's optimum]
For a single pair, find the margin $h^*$ that minimises
$-(1-\varepsilon)\log\sigma(h)-\varepsilon\log\sigma(-h)$, and the minimum value. Interpret
$\varepsilon\to0$.
\end{exercise}

\begin{exercise}[IPO vs DPO on deterministic data]
Consider one prompt with two responses where $y_w$ is \emph{always} preferred. Parameterise the
policy by a single logit $z$, with $\pi(y_w)=\sigma(z)$ and $\piref$ uniform. Write the DPO and IPO
losses as functions of $z$ and find their minimisers. Which one keeps $\pi$ away from the
deterministic policy?
\end{exercise}

\begin{exercise}[Sum vs mean]
Show that replacing summed with mean log-probs in DPO is equivalent to DPO with a
per-response $\beta_y=\beta/|y|$. For a winner twice as long as the loser, which side does this
favour, compared with standard DPO?
\end{exercise}

\begin{exercise}[Step-0 sanity check, by hand]
In Lab 26, deliberately build the reference cache with a different truncation length than training
uses. Predict, then measure, the step-0 loss and reward accuracy. Repeat with dropout left on in
the policy forward. Write a three-line assertion block for your training script that catches
both bugs.
\end{exercise}

\begin{exercise}[Prompt sharing]
Implement the packed ``one prompt, two continuations'' layout of the second Systems Lens for the toy
model (Qwen3 uses RoPE, so give $y_l$ positions starting right after the prompt). Verify that the
log-probs match \ident{sequence_logps} on the pair batch, and measure tokens processed per pair
for prompt lengths 4, 16, 64.
\end{exercise}

\begin{exercise}[Logit memory]
For Qwen3-0.6B ($V=151{,}936$), compute the fp32 logit memory for a DPO micro-batch of $B=4$ pairs
at $T=2048$. How much is saved by (a) slicing to the response half when $T_p=T_y$, (b) chunking the
LM head over 8 sequence chunks, (c) bf16 logits with fp32 \code{logsumexp}?
\end{exercise}

\begin{exercise}[Online DPO loop design]
Design a 3-round online DPO pipeline for a 7B model on 8 GPUs with a reward model of 1B parameters:
per round, generate 4 samples for each of 20k prompts, score them, form (best, worst) pairs, train one epoch.
Estimate generation time with the decode roofline of \cref{ch:inference}, training time with
the napkin math above, and decide whether generation and training should share GPUs.
\end{exercise}

% ======================================================================
\section{Lab 26: Preference optimisation from the data path up}\label{sec:dpo:lab}

\begin{labbox}{26}{Preference optimisation from the data path up}
\labdir{lab26_dpo}
\labsection{Goals}
Implement DPO and five relatives on top of a correct pair-batch data path. Show on a laptop CPU
that preference optimisation moves a measurable behaviour, and observe likelihood displacement and
the alignment tax yourself.

\labsection{Setup}
\begin{lstlisting}[style=shell]
cd code
LAB_IMPL=starter  python -m pytest labs/lab26_dpo     # your work
LAB_IMPL=solution python -m pytest labs/lab26_dpo     # reference
python labs/lab26_dpo/run.py --out results            # full experiment, minutes on CPU
\end{lstlisting}

\labsection{Tasks}
\begin{enumerate}
\item \ident{build_pair_batch}: the $(2B,T)$ layout of \cref{fig:dpo:batch} with a response mask.
\item \ident{sequence_logps}: shifted, masked sums using \code{gather} $-$ \code{logsumexp}; also return
  token counts.
\item \ident{dpo_loss} with conservative label smoothing; return detached implicit rewards.
\item \ident{ipo_loss}, \ident{kto_loss}, \ident{simpo_loss}, \ident{orpo_loss}.
\item \ident{precompute_reference_logps} and \ident{preference_metrics} (ties count one half).
\item \ident{dpo_train_loop}: shuffled micro-batches, gradient accumulation, clipping, per-step metrics.
\end{enumerate}
The tests check exact values, gradient directions, the step-0 $\ln 2$ property, that the cache
matches on-the-fly log-probs, that disabling LoRA reproduces the base model, and that 25 DPO
steps raise $P(\texttt{sure})$ substantially.

\labsection{Measurements}
\begin{itemize}
\item Run \texttt{run.py}: plot $P(\texttt{sure})$, greedy accuracy and the two log-probs against steps
  for $\beta\in\{0.05,0.1,0.5\}$. Expect faster behaviour change and larger log-prob drift for
  small $\beta$.
\item Compare final $P(\texttt{sure})$, accuracy and wall-clock for DPO, cDPO, IPO, KTO, SimPO, ORPO,
  CPO and LoRA-DPO. Which reference-free methods damage the reversal skill, and why?
\item Time the reference cache against one training epoch; relate the ratio to the $2N$ vs $6N$
  FLOP model.
\item On a GPU, repeat with Qwen3-0.6B and a public preference set (\ident{load_preference_jsonl}
  + \ident{tokenize_pair}); log peak memory with full DPO + reference model, full DPO + cache, and
  LoRA-DPO.
\end{itemize}

\labsection{Deliverables}
Passing tests; the plots; a table of final metrics per method; one paragraph explaining
likelihood displacement in your run using \eqref{eq:dpo:grad}.

\labsection{Stretch goals}
Add an NLL term on the chosen response (RPO-style) and show it reduces displacement; implement
online DPO on the toy task (sample from the current policy, label with the exact rule, retrain);
implement the shared-prefix packed layout and measure the savings.
\end{labbox}
